%-------------------------------------------------------------------------------
%	章节标题
%-------------------------------------------------------------------------------
\cvsection{工作经历}

%-------------------------------------------------------------------------------
%	内容
%-------------------------------------------------------------------------------
\begin{cventries}

%---------------------------------------------------------
  \cventry
    {AI全栈工程师（外派至汇丰银行）} % Job title
    {美钛技术服务（上海）有限公司广州分公司} % Organization
    {中国广州} % Location
    {2025年2月 - 至今} % Date(s)
    {
      \begin{cvitems} % Description(s) of tasks/responsibilities
        \item {作为TEKsystems外包至汇丰银行的AI工程师，利用Copilot及其API加速汇丰企业技术部财务转型平台一个模块的后端开发。该模块涵盖财务账本平台中的财务调整，管理调整和财务日记账的账目级控制。}
        \item {我在汇丰的正式头衔是Specialist Engineer，在TEKsystems的头衔是Full Stack Engineer。然而，鉴于我的工作性质以及我大量使用AI工具和Copilot技术来增强开发工作流，``AI全栈工程师''是对我角色最准确的描述。}
        \item {维护了财务数据处理（包括导入、验证和导出）的功能，同时增强了新财务表头的提交和审批工作流程。积累了会计、账本和银行系统方面的经验。}
        \item {参与了从本地开发到UAT测试和生产发布的完整开发生命周期。帮助退役运行在 WebSphere Application Server 上的遗留应用。使用 Ansible 和 Jenkins 自动化发布，协助将 Angular 升级到主版本，乐于解决各类任务——尤其是为团队解决瓶颈问题。}
        \item {负责模块的集成和API测试，使用Copilot API自动生成了约70个测试用例，涵盖Java中的Spring Filters、Python unittest、JSON截断、提示工程和区域端点。}
        \item {将自己打造成为自主的人形AI代理，通过脚本、日志、文档和分享，尽可能地实现自动化。通过LLM输出的缓存和验证，为20个定制代理提供专用提示和上下文，并在整个技术栈中使用了400个可重用的Copilot编写的脚本和1100份Copilot编写的指南，在大型银行里的工程体验开始变得良好。并且在内部以 premium requests 指标衡量的 Copilot 使用量中排名前 6\%。加入汇丰内部 AI 社区，并凭借高关注度项目 AIPlayer 获得贡献奖。}
        \item {技术栈包括Java、Spring、IBM Db2、Maven、Nexus、Angular、Python、HashiCorp Vault等技术，以及Windows、Control-M、IBM WebSphere Liberty Profile和Copilot等工具。}
        \item {与首席工程师合作开发涉及大型Excel文件处理和合并/联合工具的项目，用于比较多个银行系统之间的差距和差异，使用Python、DuckDB、React和Copilot。}
        \item {业余时间持续探索AI：在H200、RTX 4070和AMD MI300X GPU上训练小模型（nanoGPT、GPT-2 760M）并在个人知识库文本（约500万token）上微调模型；使用Qwen-27B进行推理；通过OpenRouter和Claude Code构建个人项目（过去一年消耗约20亿token）；熟练掌握PyTorch、ROCm、FlashAttention、Codex、llama.cpp、Transformers和推理技术。参加NVIDIA加速计算开发者技术交流会，向行业专家学习技术。}
      \end{cvitems}
    }

%---------------------------------------------------------
  \cventry
    {独立AI工程师} % Job title
    {个人执业} % Organization
    {中国广州} % Location
    {2023年8月 - 2025年1月} % Date(s)
    {
      \begin{cvitems} % Description(s) of tasks/responsibilities
        \item {分析并重构约30个机器学习项目的核心组件，项目来源包括PyTorch/TensorFlow教程、Coursera课程及开源仓库。获得两项Coursera认证：机器学习专项课程和深度学习专项课程。}
        \item {备考专升本课程，主攻高等数学、计算机组成原理和线性代数。通过日本音乐和TikTok视频沉浸式学习日语。}
        \item {作为全栈开发者参与AI故事机器人开发，基于Claude API实现个性化叙事功能。使用Python/Flask/React/Nginx构建提示配置系统和管理界面，部署于AWS云平台，采用Prometheus监控和ELK日志管理，ChatGPT-4辅助编码。}
        \item {维护原创技术博客（\url{https://lzwjava.github.io}，431篇），运用大语言模型实现9语种翻译、文本转语音及PDF/EPUB格式支持。集成GitHub工作流，使用LaTeX排版学术论文和简历。技术栈涵盖Python/Jekyll/Deepseek/Mistral。}
        \item {实验研究llama.cpp/嵌入模型/重排序/RAG架构及MMLU基准测试，集成Jina AI/Tavily AI Search/ElevenLabs等搜索引擎API。}
      \end{cvitems}
    }
%---------------------------------------------------------
  \cventry
    {后端工程师（外派至汇丰银行）} % Job title
    {深圳法本信息技术有限公司} % Organization
    {中国广州} % Location
    {2022年11月 - 2023年7月} % Date(s)
    {
      \begin{cvitems} % Description(s) of tasks/responsibilities
        \item {法本是中国领先的软件技术服务商。汇丰银行是全球最大的银行及金融服务机构之一。\href{https://payme.hsbc.com.hk/}{PayMe} 是汇丰面向香港居民开启的移动支付服务。}
        \item {外派至汇丰银行参与 PayMe 项目。参与自动充值功能后端开发：监控 Azure EventHub 支付事件，当用户余额不足时自动从绑定的卡片扣款。采用面向对象编程优雅处理业务逻辑，通过 AOP 实现审计日志。}
        \item {完成公司 AWS 培训后主导 API 云迁移：实现基于请求头的路由策略，确保数据库安全配置，参与微服务云原生部署。}
        \item {技术栈包括 Java/Spring/Kafka，云服务采用 Azure/Azure DevOps/AWS。}
      \end{cvitems}
    }

%---------------------------------------------------------
  \cventry
    {后端工程师（外派至星展银行）} % Job title
    {北京博彦科技股份有限公司} % Organization
    {中国广州} % Location
    {2021年12月 - 2022年11月} % Date(s)
    {
      \begin{cvitems} % Description(s) of tasks/responsibilities
        \item {博彦科技是领先的业务 IT 和咨询公司。星展银行（DBS）是东南亚资产规模最大的银行，也是亚洲最大的银行之一。外派至星展银行，参与 DBS Client Connect 和 DBS DigiBank CN 项目开发。}
        \item {在股票交易系统（DBS Client Connect）中实现股票展示、客户信息、交易前检查及订单提交功能，集成 Avaloq API 提升底层能力，运用编辑距离算法优化股票代码搜索体验。}
        \item {在数字银行（DBS DigiBank CN）项目开发共同基金、结构性理财产品等微服务。通过 PCF 日志分析生成 QPS 报告协助性能测试。开发 Karate 测试自动化工具提升测试覆盖率。}
        \item {采用 Java/Spring Cloud 技术栈，基于 PCF 云平台和 Kibana 日志系统，实践 BDD/TDD 方法论。}
      \end{cvitems}
    }

%---------------------------------------------------------
\cventry
{全栈工程师} % Job title
{个人执业} % Organization
{中国广州} % Location
{2020年1月 - 2021年11月} % Date(s)
{
  \begin{cvitems} % Description(s) of tasks/responsibilities
    \item {撰写技术博客分享知识，通过Netflix和英文原著提升语言能力，解决约500道算法题并参加Codeforces竞赛。}
    \item {实践Spark/Hadoop/Kubernetes/Docker入门教程，积累大数据和云原生技术基础。}
    \item {承接LED显示屏官网开发（lvchensign.com）、ShowMeBug企业微信集成、无纺布贸易数据爬虫、mathjax2mobi电子书工具等自由项目。}
    \item {企业微信集成项目使用Ruby on Rails实现技术面试工具对接，数据爬虫项目通过Python/Selenium自动化采集并生成商业分析报告。}
    \item {mathjax2mobi工具将含LaTeX公式的HTML转换为SVG图片，支持MOBI等电子书格式，技术栈含Python/BeautifulSoup/Selenium。}
  \end{cvitems}
}

%---------------------------------------------------------
\cventry
{创始人兼全栈工程师} % Job title
{北京平方根科技有限公司} % Organization
{中国北京} % Location
{2016年7月 - 2019年12月} % Date(s)
{
  \begin{cvitems} % Description(s) of tasks/responsibilities
    \item {在3.5年内运营了两块业务。2016年7月至2017年9月，推出并运营了知识直播平台"趣直播"。2018年1月至2019年12月，转型为软件咨询业务。}
    \item {趣直播：累计举办80场讲座，获3万用户和千万级浏览量。讲师使用 OBS 推流，用户可实时参与或查看回放。平台与微信深度集成。主力开发了3个 Web 客户端和1个服务器，代码提交量约2000次。技术栈：PHP/Vue/MySQL/Redis/LeanCloud/阿里云/SRS/微信SDK。}
    \item {软件咨询：完成50个定制项目，总营收300万元，利润约70万元。负责项目洽谈、团队管理及核心开发。主导项目包括：}
    \item {面包Live后端重构：将ThinkPHP/Node.js/Go技术栈统一迁移至Laravel，实现课程、支付、分销等模块，与喜马拉雅平台内容同步。}
    \item {江苏卫视《最强大脑》微信小程序：全栈开发益智答题游戏，通过 Redis 应对高并发场景。}
    \item {冲顶大会直播答题APP：基于 Java/Spring/Kafka 实现实时问答系统，运用 Socket.IO 处理交互，研究 SEI 时间戳同步方案保障直播流与题目同步。}
  \end{cvitems}
}

%---------------------------------------------------------
\cventry
{联合创始人兼全栈工程师} % Job title
{北京大米互娱有限公司} % Organization
{中国北京} % Location
{2015年11月 - 2016年7月} % Date(s)
{
  \begin{cvitems} % Description(s) of tasks/responsibilities
    \item {大米互娱是由包括我在内的6名互联网爱好者共同创立的公司。推出并运营了 CodeReview 平台，这是一个集代码评审、交流和分享于一体的专业平台，积累了约3000名用户。}
    \item {该平台功能包括用户管理、代码提交和评审流程、通知系统、支付集成以及活动和研讨会管理。工程师可以提交代码由专家评审以提高质量，专家则获得评审费。同时为用户提供线下研讨会和活动。}
    \item {负责后端系统及半数前端开发，技术栈包括PHP/Vue/CodeIgniter/阿里云/Ping++支付。}
  \end{cvitems}
}

%---------------------------------------------------------
\cventry
{软件工程师} % Job title
{美味书签（北京）信息技术有限公司（LeanCloud）} % Organization
{中国北京} % Location
{2014年7月 - 2015年11月} % Date(s)
{
  \begin{cvitems} % Description(s) of tasks/responsibilities
    \item {美味书签（LeanCloud）是中国领先的云计算服务商。提供包括对象存储、文件存储、Web 托管、容器、即时通讯、推送、短信和游戏后端在内的全套云服务，服务于数十万开发者用户。}
    \item {参与LeanCloud Objective-C/Java SDK开发，主导即时通讯演示应用LeanChat的iOS/Android客户端开发。}
    \item {技术栈涵盖iOS/Android SDK、Cocoapods、Xcode、Android Studio及Angular框架。}
  \end{cvitems}
}

%---------------------------------------------------------
\end{cventries}
